\section{Anatomi Bab dan Fitur Pembelajaran}
\label{sec:ch01-anatomi}

Setiap bab inti memiliki urutan yang tetap. Kotak \textbf{Outcome
Pembelajaran dan Sub-CPMK} menyatakan kompetensi yang harus dicapai.
Bagian \textbf{Materi Pokok} memberikan peta isi. Setelah itu pembaca
menemukan uraian materi, contoh, praktikum, latihan, asesmen, daftar
cek, rangkuman, dan evaluasi.

Tiga jenis kotak digunakan di seluruh buku:
\begin{itemize}[leftmargin=*]
  \item \textbf{Informasi} untuk konsep pendukung atau konteks tambahan;
  \item \textbf{Peringatan} untuk kesalahan umum dan jebakan konseptual;
  \item \textbf{Tips} untuk strategi praktis saat menulis kode atau
        mengerjakan ujian.
\end{itemize}

Kode program ditampilkan sebagai \textbf{Kode Program} dengan nomor
baris. Berkas lengkap proyek tersedia di direktori \file{code/subsetc}.
Contoh program pertama ditunjukkan pada \lstref{lst:ch01-hello}.

\includecode[
  language=C,
  caption={Program Subset-C pertama: penjumlahan dua bilangan},
  label={lst:ch01-hello}
]{../../code/subsetc/examples/hello.sc}

Program pada \lstref{lst:ch01-hello} masih sangat sederhana, tetapi
sudah mengandung seluruh unsur yang akan diolah oleh front-end:
deklarasi fungsi, deklarasi variabel, ekspresi, pemanggilan
\code{print}, dan pernyataan \code{return}.
